Persistent homology is a branch of computational algebraic topology that studies shapes and extracts features over multiple scales. In this paper, we present an unsupervised approach that uses persistent homology to study divergent behavior in agricultural point cloud data. More specifically, we build persistence diagrams from multidimensional point clouds, and use those diagrams as the basis to compare and contrast different subgroups of the population. We apply the framework to study the cold hardiness behavior of 5 leading grape cultivars, with real data from over 20 growing seasons. Our results demonstrate that persistent homology is able to effectively elucidate divergent behavior among the different cultivars; identify cultivars that exhibit variable behavior across seasons; and identify seasonal correlations.

